% =====================================================================
%  TAREA 07 - EMDS | Secciones 4.1 y 4.2  (responsable: Odin)
% =====================================================================

\documentclass[12pt,a4paper]{article}

\usepackage[T1]{fontenc}
\usepackage[utf8]{inputenc}
\usepackage{lmodern}                                  % Latin Modern
\usepackage[spanish,es-tabla,es-noshorthands]{babel} % "Tabla", sin shorthands
\usepackage[margin=2.5cm]{geometry}                   % margenes 2.5 cm
\usepackage{setspace}\onehalfspacing                  % interlineado 1.5
\usepackage{booktabs}
\usepackage{tabularx}
\usepackage{array}
\usepackage{ragged2e}
\usepackage[authoryear,round]{natbib}                 % citas APA autor-anio
\usepackage{xcolor}
\usepackage[hidelinks]{hyperref}

\setlength{\parskip}{0.4em}
\setlength{\bibsep}{0.5em}

\begin{document}

% La numeracion arranca en 4 para coincidir con el documento del equipo.
\setcounter{section}{3}

\begin{center}
{\large\bfseries Tarea 07 --- Secciones 4.1 y 4.2}\\[2pt]
{\normalsize Investigación de herramientas de prueba y análisis estático}\\[2pt]
{\footnotesize Estándares y Métricas para el Desarrollo de Software (E-EMDS-2)}
\end{center}
\vspace{0.6em}

% =====================================================================
\section{Desarrollo}

% ---------------------------------------------------------------------
\subsection{Investigación: frameworks de automatización E2E}
\label{sec:4-1}

La automatización de pruebas de extremo a extremo (\textit{end-to-end}, E2E) valida
el comportamiento del sistema desde la perspectiva del usuario final, ejercitando la
interfaz real contra una aplicación en ejecución \citep{jorgensen2022}. Para nuestro
proyecto integrador ---un wiki colaborativo construido con un frontend en
React/TypeScript y un backend en Node.js--- la elección del framework determina la
velocidad de la retroalimentación, la estabilidad de la \textit{suite} y el costo de
mantenimiento a largo plazo. A continuación comparamos las tres herramientas de
referencia del ecosistema E2E en 2026: \textbf{Selenium WebDriver}, \textbf{Cypress} y
\textbf{Playwright}.

\paragraph{Selenium WebDriver.}
Es la herramienta más madura del ecosistema y opera sobre el estándar
\textit{W3C WebDriver}, comunicándose con cada navegador a través de su \textit{driver}
correspondiente \citep{seleniumdocs}. Su mayor fortaleza es la amplitud: soporta
múltiples lenguajes (Java, Python, C\#, JavaScript, Ruby, Kotlin) y la mayor variedad de
navegadores, incluyendo motores heredados. La gestión de esperas se realiza mediante
esperas explícitas con \texttt{WebDriverWait} y condiciones esperadas, lo que evita la
fragilidad de las esperas fijas \citep{garcia2022}. Como contrapartida, exige más
código de andamiaje (\textit{boilerplate}), su curva de aprendizaje es más pronunciada y
la ejecución en paralelo depende de infraestructura adicional (Selenium Grid).

\paragraph{Cypress.}
Adopta una arquitectura distinta: ejecuta el código de prueba \emph{dentro} del propio
navegador, en el mismo bucle de eventos de JavaScript que la aplicación bajo prueba
\citep{cypressdocs}. Esto le otorga latencia de red nula, espera automática
(\textit{auto-waiting}) y un depurador con ``viaje en el tiempo'' que captura el estado
del DOM en cada paso, lo que se traduce en una experiencia de desarrollo (DX) superior.
Sus limitaciones son consecuencia de esa misma arquitectura: solo admite
JavaScript/TypeScript, su soporte de WebKit/Safari es experimental e incompleto, y la
ejecución paralela robusta requiere el servicio de pago Cypress Cloud.

\paragraph{Playwright.}
Desarrollado por Microsoft, automatiza de forma nativa los tres motores principales
---Chromium, Firefox y WebKit--- y ofrece enlaces oficiales para JavaScript/TypeScript,
Python, Java y C\# \citep{playwrightdocs}. Incorpora espera automática basada en
comprobaciones de \textit{actionability} (el elemento debe ser visible, estable, habilitado
y no estar obstruido antes de interactuar), contextos de navegador aislados, un
\textit{Trace Viewer} para el análisis post-mortem y, de forma destacada, ejecución
paralela y \textit{sharding} nativos sin necesidad de un \textit{grid}. Estas
características reducen la cantidad de pruebas inestables (\textit{flaky}) y disminuyen el
costo total de propiedad al escalar la \textit{suite}.

\paragraph{Tabla comparativa.}
La Tabla~\ref{tab:e2e} sintetiza la comparación según diez criterios técnicos relevantes
para nuestro contexto.

\begin{table}[htbp]
\centering
\footnotesize
\caption{Comparación de frameworks de automatización E2E (2026).}
\label{tab:e2e}
\renewcommand{\arraystretch}{1.25}
\begin{tabularx}{\textwidth}{>{\RaggedRight\arraybackslash}p{0.18\textwidth}
                             >{\RaggedRight\arraybackslash}X
                             >{\RaggedRight\arraybackslash}X
                             >{\RaggedRight\arraybackslash}X}
\toprule
\textbf{Criterio} & \textbf{Selenium} & \textbf{Cypress} & \textbf{Playwright} \\
\midrule
Lenguajes soportados &
Java, Python, C\#, JavaScript, Ruby, Kotlin &
Solo JavaScript / TypeScript &
JavaScript/TypeScript, Python, Java, C\# \\

Arquitectura &
Protocolo W3C WebDriver vía \textit{drivers} externos &
Ejecuta dentro del navegador (mismo bucle de JS) &
Comunicación directa con el navegador (WebSocket/CDP) \\

Curva de aprendizaje &
Pronunciada; requiere más \textit{boilerplate} &
Baja para equipos de JS; DX muy amigable &
Media; API moderna y \textit{codegen} \\

Velocidad / rendimiento &
Moderada; sobrecarga de WebDriver &
Buena en local; arranque más lento &
Alta; la más rápida en ejecución paralela \\

Manejo de esperas &
Explícitas (\texttt{WebDriverWait}); evita esperas fijas &
Espera automática integrada &
Espera automática por \textit{actionability} \\

Soporte de navegadores &
El más amplio (incluye motores heredados) &
Chromium y Firefox nativos; WebKit experimental &
Chromium, Firefox y WebKit nativos \\

Ejecución en paralelo &
Mediante Selenium Grid (infraestructura adicional) &
Limitada; paralelismo robusto requiere Cypress Cloud &
Nativa, sin \textit{grid} (\textit{workers} y \textit{sharding}) \\

Pruebas móviles &
Nativas vía Appium &
No (solo emulación de \textit{viewport}) &
No (solo emulación de dispositivos) \\

Comunidad / madurez &
La más madura; estándar de facto empresarial &
Fuerte en \textit{frontend} JS; comunidad consolidada &
Crecimiento más rápido; comunidad en expansión \\

Licencia &
Apache 2.0 (código abierto) &
MIT (código abierto; \textit{Cloud} de pago) &
Apache 2.0 (código abierto) \\
\bottomrule
\end{tabularx}
\\[2pt]
{\footnotesize Fuente: elaboración propia con base en la documentación oficial de
Selenium, Cypress y Playwright \citep{seleniumdocs, cypressdocs, playwrightdocs} y
en \citet{garcia2022}.}
\end{table}

\paragraph{Recomendación justificada.}
Para el contexto de nuestro proyecto ---una aplicación web de página única (SPA) con uso
intensivo de JavaScript y un \textit{stack} de TypeScript/React/Node.js--- la
recomendación técnica es \textbf{Playwright}. Tres razones lo respaldan: (1)~sus enlaces
de primera clase para TypeScript se alinean directamente con el lenguaje del proyecto,
reduciendo la fricción de adopción; (2)~su soporte nativo de WebKit permite cubrir Safari,
navegador relevante para un wiki de acceso público que Cypress no cubre de forma fiable; y
(3)~su ejecución paralela nativa, junto con la espera por \textit{actionability}, ataca
directamente el principal problema de las \textit{suites} E2E ---la inestabilidad--- sin
costos de infraestructura adicionales, en línea con el principio de pruebas rápidas y
confiables que recomienda \citet{aniche2022}.

No obstante, conforme a lo trabajado en la Semana 7, la \emph{implementación} práctica
(sección~4.3 y siguientes) se realiza con \textbf{Selenium + pytest} aplicando el
patrón \textit{Page Object Model}. Esta diferencia entre la herramienta recomendada y la
implementada es deliberada y está prevista por la asignatura: Selenium ofrece el soporte
más amplio de lenguajes y navegadores, una base conceptual sólida sobre el estándar
WebDriver y un ecosistema empresarial maduro, lo que lo hace idóneo como herramienta de
aprendizaje \citep{garcia2022}. Las técnicas de diseño de casos (válidos, de frontera
y de error) que aplicamos son independientes del framework \citep{jorgensen2022}, por
lo que la suite construida con Selenium puede migrarse a Playwright en el futuro
conservando su valor.

% ---------------------------------------------------------------------
\subsection{Investigación: análisis estático de código}
\label{sec:4-2}

El análisis estático examina el código fuente \emph{sin ejecutarlo}, identificando
defectos a partir de su estructura sintáctica y semántica. Es una práctica de
``desplazamiento a la izquierda'' (\textit{shift-left}) que detecta problemas en las
etapas tempranas del ciclo de vida, cuando corregirlos es más barato. En esta sección
comparamos tres herramientas representativas de categorías distintas
---\textbf{SonarQube}, \textbf{Semgrep} y los \textit{linters} \textbf{ESLint/Pylint}---
y diferenciamos el análisis estático del dinámico.

\paragraph{SonarQube.}
Es una plataforma de calidad y seguridad de código que cubre más de 30 lenguajes mediante
miles de reglas predefinidas \citep{sonarqubedocs}. Detecta \textit{bugs}
(confiabilidad), vulnerabilidades y \textit{security hotspots} (seguridad) y \textit{code
smells} (mantenibilidad), y además mide cobertura, duplicación y complejidad ciclomática y
cognitiva. Su característica distintiva son las \textit{Quality Gates}: criterios
automáticos de aprobación o rechazo que pueden bloquear la fusión (\textit{merge}) de una
rama cuando el código nuevo no cumple los umbrales definidos. Es importante una precisión
para la sección 4.4: bajo el modelo \textit{Clean Code} actual, SonarQube define
\emph{tres} cualidades de software ---Seguridad, Confiabilidad y Mantenibilidad--- que,
sumadas a las medidas de \emph{Cobertura} y \emph{Duplicación} del tablero, conforman las
cinco dimensiones que reporta la herramienta \citep{sonarqubedocs}.

\paragraph{Semgrep.}
Es un motor de análisis estático orientado a la seguridad (SAST) que opera mediante
coincidencia de patrones sobre el árbol de sintaxis abstracta, no sobre texto plano. Su
fortaleza es la rapidez (escaneos del orden de segundos) y la facilidad para escribir
reglas personalizadas en YAML, lo que lo hace popular para hacer cumplir políticas de
seguridad específicas. Detecta patrones explotables ---inyección SQL, \textit{cross-site
scripting} (XSS), credenciales codificadas, deserialización insegura--- mediante análisis
de contaminación (\textit{taint analysis}). Su limitación es que está enfocado en
seguridad: no calcula métricas de mantenibilidad, complejidad, duplicación ni cobertura.

\paragraph{ESLint / Pylint (linters).}
Representan la capa más rápida y cercana al desarrollador. ESLint analiza JavaScript y
TypeScript mediante un sistema de reglas configurable basado en el AST, mientras que Pylint
hace lo equivalente para Python. Aplican estilo y convenciones, detectan errores comunes y
se ejecutan en tiempo real dentro del IDE. Su limitación es la profundidad: no realizan
análisis de seguridad entre archivos ni detectan errores lógicos o de tiempo de ejecución;
por ello suelen complementarse con una herramienta SAST o una plataforma de calidad.

\paragraph{Tabla comparativa.}
La Tabla~\ref{tab:static} resume las diferencias. Se incluye CodeClimate (hoy en
transición a ``qlty'') como referencia adicional de plataforma orientada a la
mantenibilidad.

\begin{table}[htbp]
\centering
\footnotesize
\caption{Comparación de herramientas de análisis estático (2026).}
\label{tab:static}
\renewcommand{\arraystretch}{1.25}
\begin{tabularx}{\textwidth}{>{\RaggedRight\arraybackslash}p{0.16\textwidth}
                             >{\RaggedRight\arraybackslash}X
                             >{\RaggedRight\arraybackslash}X
                             >{\RaggedRight\arraybackslash}X}
\toprule
\textbf{Criterio} & \textbf{SonarQube} & \textbf{Semgrep} & \textbf{ESLint / Pylint} \\
\midrule
Categoría &
Plataforma de calidad + seguridad &
Motor SAST (seguridad) &
\textit{Linter} (estilo + errores básicos) \\

Defectos que detecta &
\textit{Bugs}, vulnerabilidades, \textit{code smells}, duplicación &
Vulnerabilidades de seguridad y \textit{anti-patrones} &
Estilo, convenciones y errores comunes \\

Enfoque de detección &
Miles de reglas predefinidas por lenguaje &
Coincidencia de patrones sobre el AST (reglas en YAML) &
Reglas configurables sobre el AST \\

Métricas de calidad &
Sí: complejidad, duplicación, cobertura, deuda técnica &
No (solo seguridad) &
No (solo reglas de estilo/sintaxis) \\

Integración CI/CD &
\textit{Quality Gates} que bloquean \textit{merges} &
Rápida; ideal para escaneo en cada \textit{pull request} &
En cada \textit{commit}; muy rápida \\

Lenguajes &
30+ lenguajes &
35+ lenguajes &
ESLint: JS/TS; Pylint: Python \\

Licencia / costo &
\textit{Community Edition} gratuita; ediciones de pago &
Núcleo de código abierto (LGPL); gratis hasta cierto número de contribuyentes &
Código abierto (gratuito) \\
\bottomrule
\end{tabularx}
\\[2pt]
{\footnotesize Fuente: elaboración propia con base en la documentación oficial de las
herramientas \citep{sonarqubedocs}.}
\end{table}

\paragraph{Análisis estático frente a análisis dinámico.}
La diferencia fundamental es la \emph{ejecución}. El análisis estático inspecciona el
código sin correrlo y, por ello, puede razonar sobre \emph{todas} las rutas del programa,
detectando \textit{code smells}, vulnerabilidades por patrón, complejidad excesiva y
duplicación antes de que el software se ejecute siquiera. Su contrapartida son los falsos
positivos y la imposibilidad de observar comportamientos que solo emergen en tiempo de
ejecución. El análisis dinámico, en cambio, examina el programa \emph{mientras se ejecuta}
con entradas concretas: detecta fallos reales, fugas de memoria, cuellos de botella de
rendimiento y condiciones de carrera, pero solo sobre las rutas efectivamente ejercitadas
(de ahí su riesgo de falsos negativos) y requiere un entorno de ejecución y datos de
prueba. De hecho, la propia \textit{suite} E2E con Selenium de la sección~4.3 constituye
una forma de análisis dinámico, pues valida la aplicación en funcionamiento. Ambos enfoques
son complementarios y se integran en el Plan de Pruebas de la Unidad~II: SonarQube (estático)
y la suite de pruebas con pytest/Selenium (dinámico) cubren clases de defectos distintas
\citep{aniche2022, khorikov2020}.

\paragraph{Justificación para nuestro \textit{stack}.}
Recomendamos \textbf{SonarQube} como herramienta principal de análisis estático para el
proyecto por cuatro motivos: (1)~analiza de forma nativa JavaScript y TypeScript, los
lenguajes de nuestro frontend y backend; (2)~unifica en un solo tablero las cinco
dimensiones que pide la evaluación (Confiabilidad, Seguridad, Mantenibilidad, Cobertura y
Duplicación), facilitando la sección~4.4; (3)~sus \textit{Quality Gates} se integran con el
pipeline de CI de la sección~4.5 para bloquear fusiones que no cumplan los umbrales; y
(4)~su \textit{Community Edition} es gratuita y, mediante SonarCloud, el análisis es
gratuito para repositorios públicos. Como complemento, conservamos \textbf{ESLint} ---ya
estándar en proyectos de React/Node.js--- como capa rápida de \textit{linting} en cada
\textit{commit}, y dejamos abierta la incorporación de \textbf{Semgrep} si el equipo
requiere profundizar en seguridad. Esta estratificación (\textit{linter} en cada commit,
plataforma de calidad en cada PR) es la práctica recomendada para mantener la salud del
código sin sacrificar velocidad \citep{aniche2022}.

% =====================================================================
%  REFERENCIAS  (APA 7, incrustadas; natbib autor-anio)
% =====================================================================
\begin{thebibliography}{99}
\setlength{\itemsep}{0.4em}

\bibitem[Aniche(2022)]{aniche2022}
Aniche, M. (2022). \textit{Effective software testing: A developer's guide}. Manning Publications.

\bibitem[Cypress.io(2026)]{cypressdocs}
Cypress.io. (2026). \textit{Cypress documentation}. \url{https://docs.cypress.io/}

\bibitem[García(2022)]{garcia2022}
García, B. (2022). \textit{Hands-on Selenium WebDriver with Java}. O'Reilly Media.

\bibitem[ISTQB(2023)]{istqb2023}
International Software Testing Qualifications Board. (2023). \textit{Certified tester foundation level syllabus v4.0}. ISTQB. \url{https://www.istqb.org/}

\bibitem[Jorgensen(2022)]{jorgensen2022}
Jorgensen, P. C. (2022). \textit{Software testing: A craftsman's approach} (5.\textsuperscript{a} ed.). CRC Press.

\bibitem[Khorikov(2020)]{khorikov2020}
Khorikov, V. (2020). \textit{Unit testing principles, practices, and patterns}. Manning Publications.

\bibitem[Microsoft(2026)]{playwrightdocs}
Microsoft. (2026). \textit{Playwright documentation}. \url{https://playwright.dev/}

\bibitem[Selenium Project(2026)]{seleniumdocs}
Selenium Project. (2026). \textit{Selenium documentation}. \url{https://www.selenium.dev/documentation/}

\bibitem[SonarSource(2026)]{sonarqubedocs}
SonarSource. (2026). \textit{SonarQube documentation}. \url{https://docs.sonarsource.com/}

\end{thebibliography}

\end{document}
